\section{Introduction}

Repeated uses of an ordered binary quantum measurement resolve
different perturbations at different scales.  An interior change of
an eigenvalue is detected at scale $N^{-1/2}$, whereas a rotation of
a projective measurement can be detected at scale $N^{-1}$.  These
scales meet at singular effects and repeated eigenvalues.  We
determine the resulting geometry on the compact body of matrix
effects, calculate its description entropy, and construct common
learners whose loss is distinguishability by all future experiments
with at most $N$ uses.

For an input dimension $d$, let
\[
 \mathfrak E_d=\{E=E^*\in M_d(\C):0\preceq E\preceq I_d\}.
\]
An effect $E$ specifies the memoryless ordered binary measurement
\[
 \mathcal M_E(\rho)=\operatorname{tr}(E\rho)|1\rangle\langle1|
       +\operatorname{tr}((I_d-E)\rho)|0\rangle\langle0|.
\]
The input is consumed; the device returns a classical outcome and
no residual quantum system.  The distance $d_N$ is the supremum
of the unhalved trace norm between the final states of a common
experiment with at most $N$ calls.  Such an experiment may entangle
inputs, retain an arbitrary reference and quantum memory, use
feedback, and stop at a public bounded time.  Thus $0\le d_N\le2$;
classical total variation is one half of the corresponding trace
norm.  The nonadaptive restriction $d_N^{\rm na}$ allows entangled
input blocks and retained references, without feedback between calls.

\subsection{Geometry and entropy of the closed effect body}
For $E,F\in\mathfrak E_d$, put
\[
 M=\tfrac12(E+F),\qquad H=F-E,\qquad V=M(I_d-M),
\]
and define the comparison modulus
\[
 Q_N(E,F)^2=N\langle H,
       (L_V+R_V+N^{-1}\operatorname{Id})^{-1}H\rangle_{\rm HS}.
\]
The inverse acts on the real Hilbert space of Hermitian matrices.
For every $d,N\ge1$, Theorem~\ref{thm:matrixmetric76} gives
\begin{equation}\label{eq:intrometric77}
 \frac{\min\{1,Q_N(E,F)\}}{8192d}
 \le d_N^{\rm na}(\mathcal M_E,\mathcal M_F)
 \le d_N(\mathcal M_E,\mathcal M_F)
 \le\min\{2,8Q_N(E,F)\}.
\end{equation}
The comparison holds at every spectral rank and multiplicity, with
no margin from the support endpoints.  In particular,
$d_N\le65536d\,d_N^{\rm na}$ uniformly in $N$.  The modulus
provides an explicit comparison of the operational distances;
its use requires no triangle inequality for $Q_N$.

A horizontal measurement dilation gives the adaptive tangent bound.
A regularized Sylvester equation separates the tangent into a
horizontal part and a remainder controlled by the hybrid inequality.
Operator concavity of $G(I_d-G)$ then integrates the bound along
a finite segment, including boundary endpoints.  Bernoulli inputs
and exact two-dimensional compressions give the converse.  These
pair-dependent experiments establish separation; the common
learning construction below is a separate argument.

Let $\mathcal C_N(\mathfrak E_d,\delta)$ be the least number of
legal memoryless binary centres in a $d_N$-cover.  For every fixed
$d$, Theorem~\ref{thm:matrixcover76} proves
\begin{equation}\label{eq:introentropy77}
 \mathcal C_N(\mathfrak E_d,\delta)
 \asymp_d N^{d^2/2}[\log(N+2)]^{\lfloor d/2\rfloor}\delta^{-d^2},
 \qquad N\ge1,\quad0<\delta\le\delta_d.
\end{equation}
The positive cap $\delta_d$ and comparison constants may depend
on $d$.  The proof bounds the weighted volume of actual operational
balls uniformly over all boundary strata.  The total volume is a
spectral integral: its ordered endpoint sectors have cumulative
exponents $S_j^2/2$, where $S_j$ is a signed prefix sum.  Balanced
prefixes, together with matching nested endpoint boxes, give the
logarithmic multiplicity $\lfloor d/2\rfloor$.

\subsection{Common learning and independent acquisition blocks}
The learner receives calls to an unknown effect, without its matrix
entries.  Theorem~\ref{thm:matrixlearning77} constructs one procedure
common to all $E\in\mathfrak E_d$ such that
\begin{equation}\label{eq:introlearning77}
 \sup_E\mathbb P_E\{d_N(\mathcal M_E,\mathcal M_{\widehat E})
                         >\delta\}\le\eta,
 \qquad
 M\le Cd^4N\delta^{-2}\log(d/\eta),
\end{equation}
where $C$ is absolute, $N\ge1$, $0<\delta\le\delta_0$ for an
absolute $\delta_0>0$, and $0<\eta\le1/8$.  The training-call
bound holds on every record, and each entangled block uses at most
$N$ calls.  A scalar Bernoulli subfamily gives a lower bound of
order $N\delta^{-2}\log(1/\eta)$ even for arbitrary
reference-assisted adaptive training.  Thus the horizon, accuracy,
and confidence orders are optimal for fixed $d$; the displayed
$d^4$ factor is an upper bound.

The qubit construction learns unknown bias, contrast, and direction.
Independent fair signs remove the bias from GHZ parity observations.
A dyadic amplitude search maintains a small-cone invariant and
terminates at the first rejected scale.  Fresh observations after
each data-dependent choice complete the direction and endpoint
estimates.  Conditional confidence allocation controls exceptional
records as well as the maximum training cost.

The matrix extension first calibrates both support boundaries in
relative Loewner order at scale $N^{-1}$ and the absolute matrix
error at scale $N^{-1/2}$.  In the calibrated eigenbasis, these
estimates control the additional variance produced by every
two-dimensional compression.  Qubit learning on the compressions
then supplies compatible coordinates in the regularized quadratic
form.  Projection onto the legal effect body, followed by rational
selection, completes the estimate.  Conditional summation of the
calibration and compression budgets introduces no $\log N$ factor.

The maximal independent block length is a further acquisition
resource.  Within a fresh block, reference-assisted quantum controls
are allowed.  At its boundary, completed outputs may be processed
jointly, and the classical record may determine the next preparation.
The next block, including its reference and unused inputs, is
isolated from earlier quantum registers until its final call.
Thus collective final estimation and classical feedback are allowed,
while older quantum memory cannot feed unfinished queries.

For a declared cap $1\le b\le N$ and every fixed $d\ge2$,
Theorem~\ref{thm:blocklearning78} proves
\begin{equation}\label{eq:introblock78}
 M_b^\star(d,N,\delta,\eta)
 =\Theta_d\!\left(\frac{N^2}{b\delta^2}\log\frac1\eta\right)
\end{equation}
for an absolute small-error range and $0<\eta\le1/8$.  The
explicit upper bound is
\[
 Cd^4(N^2/b)\delta^{-2}\log(d/\eta).
\]
Horizon blocking gives the common learner; a history-weighted
root-fidelity bound gives the converse with arbitrary storage and
collective processing of completed blocks.  In dimension one the
device is a coin, and the cost is
$\Theta(N\delta^{-2}\log(1/\eta))$ for every $b$.

\subsection{Joint learning laws in the interior}
On the full-dimensional interior
\[
 \mathfrak I_d=\{E:I_d/4\preceq E\preceq3I_d/4\},
\]
the operational distance has a dimension-free comparison with
$\min\{1,\sqrt N\|E-F\|_{\rm op}\}$.  The binary Choi estimator
of Mele--Bittel \cite{MB67}, analysed under this loss, gives
\begin{equation}\label{eq:introinterior78}
 M_{\rm int}^\star(d,N,\delta,1/8)
 =\Theta(d^2N\delta^{-2})
\end{equation}
with absolute constants, using independent one-call probes.  For
the converse, a family within operator radius $r$ of $I/2$
supplies at most $4r^2$ information about its finite unknown label
per call.  Matrix packing and Fano's inequality give
Theorem~\ref{thm:dimensionlower78}, which permits arbitrary coherent
adaptive training.

The confidence dependence is also determined jointly.  For all
$d,N\ge1$, $0<\delta\le2^{-13}$ and $0<\eta\le1/8$,
Theorem~\ref{thm:interiorconfidence80} proves
\begin{equation}\label{eq:introconfidence80}
 M_{\rm int}^{\star}(d,N,\delta,\eta)
 =\Theta\!\left(\frac{N}{\delta^2}
                   \left(d^2+d\log\frac1\eta\right)\right).
\end{equation}
The constants are absolute.  The upper bound uses independent
one-call Choi acquisitions, and the lower bound allows coherent
adaptive training; hence the order is the same under every
fresh-block cap.  At $N=1$ it gives the joint dimension, accuracy,
and confidence law for binary operator-norm learning, including
the fixed-confidence order $\Theta(d^2\epsilon^{-2})$.

The confidence converse compares $I_d/2$ with $d$ diagonal
single-coordinate perturbations.  Revealing the sampled basis
index gives a relative-entropy chain rule through retained quantum
memory.  Under the central hypothesis, the expected occupation
counts sum to the entire budget, so binary testing forces the
$d\log(1/\eta)$ term.  Taking the maximum with the packing
bound gives the sum in \eqref{eq:introconfidence80} up to an
absolute factor.  For the upper bound, the estimator of
\cite[Corollary~III.9]{MB67} is run at a rational failure
probability exponentially small in $d$.  Independent repetition
and metric-majority selection give the required confidence at the
same additive cost.

The resulting operator estimator also applies to the closed body.
At operator accuracy $\delta/(2N)$, the hybrid inequality gives
future loss at most $\delta$.  Selection between this one-call
procedure and the block learner, using only public parameters,
yields
\[
 M_b^\star\le C\frac{N^2}{\delta^2}
 \min\left\{d^2+d\log(1/\eta),\frac{d^4}{b}\log(d/\eta)\right\}
\]
in their common range; see
\eqref{eq:fullbodyconfidenceupper80}.  The corresponding joint
lower obstructions are given in
\eqref{eq:fullbodyconfidencelower80}.  In particular, the projective
boundary in \eqref{eq:introblock78} accounts for the different
horizon cost of one-call acquisition on the closed body.

\subsection{Exact descriptions and learned words}
Theorem~\ref{thm:matrixcodec77} constructs a deterministic finite
rational codebook at the order in \eqref{eq:introentropy77}.
Exact comparisons of $Q_N^2$ define a greedy selection on a finite
legal matrix grid.  Separation in the actual operational metric
controls its cardinality, and inward rounding gives coverage at
every rank.  Encoder and decoder reconstruct the same dictionary
from public parameters; every fixed-length word decodes to a legal
effect.  Its payload is
\[
 \frac{d^2}{2}\log_2N+
 \lfloor d/2\rfloor\log_2\log(N+2)
             +d^2\log_2(1/\delta)+O_d(1).
\]
The direct qubit code also gives detailed resource bounds.
Corollary~\ref{cor:matrixlearnedcode77} combines the common learner
and exact encoder at their respective optimal fixed-dimensional
orders, without further device calls.
Corollary~\ref{cor:blocklearnedcode78} gives the same conclusion
under a prescribed independent-block cap.

The interior has a dimension-uniform description law.  For all
$d,N\ge1$ and $0<\delta\le2^{-13}$,
Theorem~\ref{thm:interiorentropy79} gives
\begin{equation}\label{eq:introinteriorentropy79}
 \left(\frac{\sqrt N}{2048\delta}\right)^{d^2}
 \le\mathcal C_N(\mathfrak I_d,\delta)
 \le\left(\frac{25\sqrt N}{\delta}\right)^{d^2}.
\end{equation}
The lower bound allows arbitrary legal binary centres.
Theorem~\ref{thm:interiorcodec79} realizes the upper bound by
finite exact rational comparison and public reconstruction.  With
\[
 A_d(N,\delta)=\tfrac12d^2\log_2N+d^2\log_2(1/\delta),
\]
its length is $A_d(N,\delta)+O(d^2)$, uniformly in all parameters;
there is no $d^2\log d$ entrywise-transmission penalty.  The
absence of the endpoint logarithm reflects the interior geometry.

With shared randomness and worst-target failure probability
$0\le\eta\le1/8$, Theorem~\ref{thm:prefixrate79} gives optimal
expected one-way prefix length
$(1-\eta)A_d(N,\delta)+O(d^2)$.  Its converse permits a
randomized decoder and counts all nonpublic information.  Its
upper bound uses a public target-independent abort event, whose
occurrence requires no transmitted flag.  For learned descriptions,
Theorem~\ref{thm:interiorlearnedcode79} attains the call order
in \eqref{eq:introinterior78} and this deterministic index length
at failure $1/8$; Corollary~\ref{cor:interiorconfidencecode80}
attains \eqref{eq:introconfidence80} at every stated confidence.

\subsection{Finite risk certificates and trusted controls}
For rational public errors, these joint query and description
orders admit a complete finite classical specification before the
first device call.  The algebraic construction in
Theorem~\ref{thm:effectiveinterior80} fixes the public dictionary
and a collective readout on independent Choi records.  The
subnormalized record is polynomial in the unknown effect, so
positivity, completeness, and uniform success form a first-order
formula over the reals with rational coefficients.  Effective
quantifier elimination decides feasibility, and algebraic-point
selection returns exact readout data.  The statistical existence
bound controls the first feasible call count.

Section~\ref{sec:finiterisk81} gives a rational construction with
finite certificates of uniform risk.  Theorem~\ref{thm:finitecertificate81}
transfers verified inequalities on a finite rational grid to every
real effect in $\mathfrak I_d$.  If the grid is an operator-norm
$r$-net, and its tests bound failure by $\alpha$ at operator
radius $a$, then an $m$-call readout satisfies
\[
 \sup_{E\in\mathfrak I_d}
   \mathbb P_E\{\|E-C_J\|_{\rm op}>a+r\}\le\alpha+mr.
\]
An analytic continuity bound for the complete Choi record proves
the transfer with one fixed set of good labels at each grid point.
The prescribed parameters certify operational loss $5\delta/8$
with failure at most $5\eta/16$.
Proposition~\ref{prop:rationalreadout81} supplies rational POVMs
with exact positivity and completeness while preserving the risk
margin.  Bounded-denominator enumeration and a doubling schedule
then give the computable learner of
Theorem~\ref{thm:rationallearner81}, with the call order in
\eqref{eq:introconfidence80} and payload
$A_d(N,\delta)+O(d^2)$.  This construction uses finite rational
tests in place of real quantifier elimination.

The finite-control theorem approximates trusted preparations and
readouts with a pathwise norm budget.  It preserves the operational
loss, confidence, and every-record call orders, and it charges
classical specifications and trusted realization errors explicitly.
The interior constructions retain fresh one-call acquisitions and
collective processing of completed outputs.  Exhaustive dictionary
reconstruction, readout synthesis, quantum storage, and known-control
work are separate resources from device calls and transmitted bits;
the finite constructions assert computability without an efficiency
bound for those resources.

\subsection{Antecedents and organization}
Horizontal Kraus gauges, channel-extension metrology, and integration
of local distinguishability underlie the path argument
\cite{FI76,DKG67,DM76,KGAD76,YF76,SD76}.  Inverse-Sylvester
geometry and spectral matrix integration have classical antecedents
\cite{Dittmann76,SZGeom76}.  Here these methods enter the explicit
closed-body midpoint comparison and the nested endpoint entropy.
Mele--Bittel \cite{MB67} and
Zambrano--Ramos-Calderer--Kueng \cite{ZRK76} provide common
estimators under one-use measurement losses.  Their precise access,
dimension, rank, and confidence statements, and the future-loss
analysis of those acquisitions, are compared in
Section~\ref{sec:matrixcomparison76}.  The phase-estimation
components of the full-body learner use ideas from
\cite{HB76,KLY76}, with the conditional construction proved here.

The matrix geometry and entropy precede the qubit learning
construction, its matrix extension, and the exact dictionary.
The independent-block theorem and the interior comparison then
lead to the joint learning, confidence, and description laws.
Finite-control synthesis and rational risk certification complete
the construction.  The subsequent qubit geometry and the final
prerequisite appendix provide the remaining proofs needed by the
article.
